\section{From Vibe Coding to Verified Merge}

An agent opens a pull request that changes a checkout flow. The diff touches product UI, a generated API client, a schema-adjacent file, tests, and documentation. The hard reviewer question is no longer whether the patch looks plausible. It is: what evidence makes this safe to merge?

That question is the center of Jankurai. AI-assisted coding makes plausible code cheap; it makes ambiguous repositories expensive. GitHub's 2025 Octoverse report describes AI as mainstream in development and TypeScript as the most-used language on GitHub \cite{githubOctoverse2025}. DORA's 2025 research frames AI as an amplifier of existing organizational capability rather than an automatic productivity guarantee \cite{doraStateAIAssistedSoftwareDevelopment2025}. Security evidence is less forgiving: Veracode reports syntactically plausible generated code with known flaws \cite{veracodeGenaiCodeSecurity2025}, and GitGuardian reports accelerating secret exposure in public GitHub activity, including AI-assisted commits \cite{gitguardianSecretsSprawl2026,gitguardianSecretsSprawl2026Blog}.

The correct response is not to slow agents down until the old review process survives. The correct response is to redesign the repository as the alignment layer. The repository, not the prompt or the reviewer memory, must encode ownership, allowed edit scope, proof lanes, generated boundaries, permission limits, evidence receipts, repair queues, and versioned conformance.

Jankurai's operating claims are deliberately concrete:
\begin{itemize}
\item The repository is the alignment layer for human- and agent-authored code.
\item Vibe artifacts are detectable as repository states, not merely as reviewer discomfort.
\item Vibe artifacts are reducible when findings route to owners, proof lanes, receipts, and repair queues.
\item Reducing vibe artifacts improves agent efficiency by replacing broad context rereads and improvised proof selection with path-scoped policy and reusable evidence.
\end{itemize}

\begin{figure*}[t]
\centering
\fbox{\begin{minipage}{0.94\textwidth}
\centering
\small
change proposal
\(\rightarrow\) changed paths
\(\rightarrow\) vibe artifact scan
\(\rightarrow\) owner route
\(\rightarrow\) proof lane
\(\rightarrow\) receipt
\(\rightarrow\) merge witness
\(\rightarrow\) \texttt{pass}/\texttt{review}/\texttt{block}/\texttt{ratchet\_fail}/\texttt{release\_fail}
\(\rightarrow\) repair queue or waiver
\(\rightarrow\) reduced baseline
\end{minipage}}
\caption{The Jankurai continuous repository alignment loop. A change is not trusted because it looks plausible; it is trusted only when changed paths, owners, proof receipts, missing evidence, artifact digests, and commit/tool identity are bound into a merge witness.}
\label{fig:alignment-loop}
\end{figure*}

The checkout-flow pull request is the running example. If the generated client changed by hand, the repository should emit \texttt{HLT-002-GENERATED-MUTATION}. If the UI changed without screenshot, ARIA, trace, or geometry evidence, it should emit \texttt{HLT-013-RENDERED-UX-GAP}. If the path has no owner or proof route, it should emit \texttt{HLT-003-OWNERLESS-PATH} or \texttt{HLT-004-UNMAPPED-PROOF}. The merge witness then records whether the PR can pass, needs human review, must block, violates a ratchet, or fails release gates.

Agent-native engineering is not ``humans approve whatever the model says.'' It is human-directed engineering in which intent becomes bounded authority, changed paths become proof lanes, proof lanes leave evidence, evidence drives repair or waiver expiry, and repeated repairs become reusable primitives. The same machinery must govern human-authored code, because a vague PR from a person and a vague patch from an agent fail in the same place: the repository cannot tell what owns the behavior, which proof applies, or what evidence would make the merge trustworthy.

This paper makes four contributions. First, it defines a stack-neutral repository conformance model for continuous repository alignment. Second, it specifies a merge-witness artifact that binds changed paths to ownership, proof, evidence, and decision. Third, it turns missing evidence into stable rule IDs and an ordered repair queue. Fourth, it reports current seed conformance-suite evidence and separates the non-normative reference profile from Jankurai Core.
